\section{What svd2rust Emits, and How to Read It}
\label{sec:svd2rust-emits}

The generated PAC looks impenetrable because it is generated: names are long, types nest three deep,
and every field has a writer type. It is worth ten minutes to learn the four shapes, because after
that the generated code is the best documentation the chip has.

\subsection{Shape 1: The Peripheral}

\begin{minted}{rust}
pub type MvpS = crate::Periph<mvp_s::RegisterBlock, 0x4d00_0000>;
\end{minted}

A peripheral is a zero-sized handle carrying its base address as a const generic. That is what makes
\codeinline{Peripherals::take()} sound: the singleton exists once, and each field is a handle with a
compile-time address.

\subsection{Shape 2: The Register Block}

\begin{minted}{rust}
#[repr(C)]
pub struct RegisterBlock {
    ipversion: Ipversion,     // 0x00
    en: En,                   // 0x04
    swrst: Swrst,             // 0x08
    cfg: Cfg,                 // 0x0c
    status: Status,           // 0x10
    ...
}
\end{minted}

The field order is the address map, and the doc comment on each accessor carries the offset — the one
piece of information you need when comparing code against a manual page.

\subsection{Shape 3: Read, Write, Modify}

Every register has three call styles, and choosing the right one is a safety decision, not a style
preference:

\begin{minted}{rust}
// read: no side effects on the hardware
let bits = mvp.status().read().bits();

// write: sets every field in the register, the others to their writer default
mvp.en().write(|w| w.en().set_bit());

// modify: read-modify-write, preserving fields you did not name
gpio.portc_model().modify(|_r, w| w.mode3().pushpull());
\end{minted}

Chapter~4's central defect — a \codeinline{write} where a \codeinline{modify} was needed, which
disabled fifteen other pins — is visible in these three signatures and invisible in the code that
calls them. That is why the rule is stated at the level of the register, not of the driver.

\subsection{Shape 4: The Field Writer}

The most useful generated code is the field writer, because it is where the SVD's enumerated values
become a Rust enum:

\begin{minted}{rust}
// From the SVD: <field><name>CMOA</name> with values None/Toggle/Clear/Set
mvp.cc0_ctrl().modify(|_r, w| w.cmoa().toggle());

// And when the SVD says a register has no such field, the compiler says so:
// error[E0599]: no method named `pwm` found for `FieldWriter<.. Cmoa, ..>`
\end{minted}

That second snippet is from this project. It is a compile error that saved a silent misconfiguration:
the timer's compare-action field genuinely has four values, and PWM is not one of them — the PWM
waveform in this part comes from the register's \emph{mode}, not from the action field. A HAL written
against a guess would have written a valid-looking bit pattern and configured the wrong waveform.

\subsection{Unsafe, and Where It Is Honest}

\begin{minted}{rust}
mvp.instr0cfg0().write(|w| unsafe { w.bits(0x0102_0000) });
\end{minted}

\codeinline{bits()} is unsafe because the type system cannot know whether the value you assembled
means what you think. That \codeinline{unsafe} is the correct amount of friction: it appears exactly
where a field-level API does not exist, and its presence in a diff is a signal that the code is at
the register level on purpose.
